\section{Exponential count moments for GOE}\label{app:goe-counts}

We prove the dimension-uniform count bound used in
Section~\ref{sec:dbm} to pass from correlation functions to Laplace
functionals.  As there, \(W_m\) is normalized GOE and
\[
 \mathcal G_m=\sum_{j=1}^m
       \delta_{(m/\pi)\lambda_j(W_m)}.
\]
The argument is independent of the Anderson model: Gaussian decimation
compares GOE counts with those of a determinantal process, whose
projection kernel gives the required exponential moments.

\begin{proof}[Proof of Lemma~\ref{lem:goe-counts}]
Fix a compact interval \(J\subset\RR\) and a number \(b>0\).
We use the Gaussian decimation identity of Forrester and Rains
\cite[arXiv version~2, Theorem~5.2 and Equation~(5.9)]{ForresterRains2001}
and the count estimate in \cite[Lemma~6.3]{bu}. We reproduce the identity
in our normalization before proving the dimension-uniform bound.
Write \(\mathrm O_m\) for the decreasingly ordered
points with density proportional to
\[
 \Delta(y)\prod_{i=1}^m b_0(y_i),\qquad
 b_0(y)=e^{-y^2/2},\qquad
 \Delta(y)=\prod_{i<j}(y_i-y_j).
\]
They have the law of the eigenvalues of \(\sqrt m\,W_m/\sqrt2\).
Thus the count \(\mathcal G_m(J)\) is the \(\mathrm O_m\) count in
\[
 J_m=\frac{\pi}{\sqrt{2m}}J,
 \qquad |J_m|=O(m^{-1/2}).
\]

Superpose independent \(\mathrm O_m\) and \(\mathrm O_{m+1}\),
and order the resulting points as
\(z_1>\cdots>z_{2m+1}\). Apart from a constant and the common Gaussian
factor, their density is the polynomial
\[
 P(z)=\sum_{\substack{S\subset\{1,\ldots,2m+1\}\\|S|=m}}
          \Delta(z_S)\Delta(z_{S^c}),
\]
where each sublist is in the inherited order. The parity--Vandermonde
factorization of
\cite[arXiv version~2, Lemma~4.1 and Equation~(4.2)]{ForresterRains2001} is
\begin{equation}
 \label{eq:dbm-decimation-polynomial}
 P(z)=c_m\Delta(z_1,z_3,\ldots,z_{2m+1})
             \Delta(z_2,z_4,\ldots,z_{2m}),\qquad c_m>0.
\end{equation}
If \(z_i=z_j\) with \(i,j\) of the same parity, terms having
both indices in one sublist vanish. The other terms cancel in pairs
under exchanging membership of \(i,j\): the sign ratio is
\((-1)^{j-i-1}=-1\), since each intervening index belongs to exactly
one sublist. Hence every factor on the right divides \(P\).
Both sides have degree \(m^2\), so the quotient is constant.
It is positive on strictly decreasing \(z\), proving
\eqref{eq:dbm-decimation-polynomial}.

Fix the even points \(y_i=z_{2i}\).
The odd points interlace them, one in each of
\[
 (y_1,\infty),\ (y_2,y_1),\ldots,(y_m,y_{m-1}),\
 (-\infty,y_m).
\]
To integrate their Vandermonde times their Gaussian factors, use
the monic polynomials
\[
 p_j(x)=(-1)^j\frac{b_0^{(j)}(x)}{b_0(x)},
 \qquad 0\le j\le m.
\]
Writing the Vandermonde as the determinant of these polynomials
reduces the integral, by multilinearity, to a determinant of
integrals over the displayed intervals. Replace each row by the
sum of it and all preceding rows. For column \(j\ge1\), the first
\(m\) rows become
\[
 \int_{y_i}^\infty p_j(x)b_0(x)\,d x
       =p_{j-1}(y_i)b_0(y_i),
\]
and the last row is zero. In column zero the last row is
\(\int_\RR b_0\).
Expansion along that last row leaves a Vandermonde in the \(y_i\)
times \(\prod_i b_0(y_i)\).
More explicitly, the ascending-degree polynomial columns give signs
\((-1)^{m(m+1)/2}\) for the odd-point Vandermonde and
\((-1)^{m(m-1)/2}\) for the even-point Vandermonde, while expansion
along the last row contributes \((-1)^m\). Their product is one.
Thus, if \(x_0>y_1>x_1>\cdots>y_m>x_m\), the integral just evaluated
is exactly
\[
 \int_{\text{interlacing}}
        \Delta(x_0,\ldots,x_m)\prod_{i=0}^m b_0(x_i)
                \,d x_0\cdots\,d x_m
     =\sqrt{2\pi}\,\Delta(y)\prod_{i=1}^m b_0(y_i).
\]
Together with the even-point factors already present in
\eqref{eq:dbm-decimation-polynomial}, this proves that the even
points have density proportional to
\begin{equation}
 \label{eq:dbm-even-density}
 \Delta(y)^2\prod_{i=1}^m e^{-y_i^2}.
\end{equation}

Let \(\psi_j(x)\) be the degree-\(j\) orthonormal polynomial for
the weight \(e^{-x^2}\), multiplied by \(e^{-x^2/2}\).
The kernel
\[
 K_m(x,y)=\sum_{j=0}^{m-1}\psi_j(x)\psi_j(y)
\]
is the rank-\(m\) orthogonal projection in \(L^2(\RR)\) onto their
span. The symmetric density in \eqref{eq:dbm-even-density} is
\((m!)^{-1}\det(K_m(y_i,y_j))_{i,j=1}^m\).
Indeed the determinant is a squared weighted Vandermonde, and
orthonormality gives its integral \(m!\).
More generally its order-\(a\) factorial moment density is
\(\det(K_m(y_i,y_j))_{i,j=1}^a\).
For completeness, expand a size-\(a\) determinant by Cauchy--Binet
as the sum of squared minors indexed by \(a\) basis functions.
Integrating its last variable uses their orthogonality, and each
minor of size \(a-1\) occurs \(m-a+1\) times. Iterating this identity
proves the factorial density formula with its stated normalization.

We need a bound uniform in both \(m\) and location.
The operator \(-\partial_x^2+x^2\) has eigenvalue \(2j+1\) on
\(\psi_j\). This follows by conjugating it by \(e^{-x^2/2}\):
on the polynomial factor it acts as
\(-p''+2xp'+p\), preserves successive degree spaces, is symmetric
for the Gaussian inner product, and has leading coefficient
multiplier \(2j+1\).
Consequently every \(f\) in the range of \(K_m\) satisfies
\[
 \|f'\|_2^2\le(2m-1)\|f\|_2^2.
\]
The elementary bound
\(|f(x)|^2\le2\|f'\|_2\|f\|_2\), obtained by integrating the
derivative of \(|f|^2\), gives
\[
 K_m(x,x)
   =\sup_{\substack{f\in\operatorname{Ran}K_m\\\|f\|_2=1}}
       |f(x)|^2
   \le2\sqrt{2m}.
\]
If \(N_m^{\mathrm{even}}\) is the number of even points in \(J_m\),
the Gram determinant inequality yields, for every \(a\ge0\),
\[
 \EE(N_m^{\mathrm{even}})_a
   \le \bigl(2\sqrt{2m}\,|J_m|\bigr)^a\le C_J^a.
\]
For \(a>m\) the left side is zero. Hence, for every \(v>0\),
\[
 \EE e^{vN_m^{\mathrm{even}}}
   =\sum_{a=0}^m
       \frac{(e^v-1)^a}{a!}\EE(N_m^{\mathrm{even}})_a
   \le\exp\bigl(C_J(e^v-1)\bigr).
\]
In any interval the total superposition count is at most twice
its even count plus one. The original \(\mathrm O_m\) count is
no greater than the superposition count. Thus
\[
 \EE e^{b\mathcal G_m(J)}
   \le e^b\,\EE e^{2bN_m^{\mathrm{even}}},
\]
which proves the assertion.
\end{proof}
